\section{Boundary examples, finite checks and scope}
\label{sec:examples}

The logarithmic family already realizes both alternatives of the theorem.
For $\beta\in\mathbb R$, let
\[
 L_\beta(x)=(\log(ex))^{-\beta},\qquad
 a_\beta(k)=\frac{1}{\sqrt{k}\,(\log(ek))^\beta}.
\]
These functions satisfy all hypotheses on $L$.
The positive function $1/(x(\log(ex))^{2\beta})$ is eventually decreasing.
Integral comparison and the substitution $t=\log(ex)$ give
\[
 \int_1^\infty\frac{dx}{x(\log(ex))^{2\beta}}
 =\int_1^\infty t^{-2\beta}\,dt.
\]
Therefore $F_\infty<\infty$ exactly when $\beta>1/2$.
At the boundary $\beta=1/2$ the sum still diverges.
For $\beta\ne0$, $a_\beta(6)\ne a_\beta(2)a_\beta(3)$ because
$(1+\log2)(1+\log3)>1+\log6$. Thus these examples include genuinely
nonmultiplicative coefficients.

\begin{table}[ht]
\centering
\small
\begin{tabularx}{\textwidth}{@{}lXX@{}}
\toprule
 & $F(N)\to\infty$ & $F(N)\to F_\infty<\infty$\\
\midrule
Entrywise limit &
Universal kernel $\gcd(m,n)/\sqrt{mn}$ &
Gram form of the maximal convolution, divided by $F_\infty$\\
Closed limiting object &
Zero strong-resolvent limit; entrywise form purely singular &
Nonzero $C^*C/F_\infty$ on its operator domain\\
First diagonal entry &
$A_N(1,1)=1$ for all $N$ &
$A_N(1,1)=1$ for all $N$\\
Logarithmic example &
$\beta\le1/2$ &
$\beta>1/2$\\
\bottomrule
\end{tabularx}
\caption{The two critical alternatives on the same space $\ell^2(\N)$
and with the same normalization $F(N)$. In the summable column,
the Gram form means $\ip{Ce_m}{Ce_n}/F_\infty$, not an assertion that
each $e_m$ belongs to $D(C^*C)$.}
\label{tab:branches}
\end{table}

\paragraph{Exact finite sentinels.}
The accompanying finite checks use rational arithmetic after the
coordinate change $f(m)=b(m)\sqrt m$, with
$b(1)=1$, $b(2)=-1$, $b(3)=1/2$, $b(6)=1/3$ and zero otherwise.
At $r_p=p^{-1/2}$ this makes
\[
 q_K[f]=\sum_{m,n}b(m)b(n)\gcd(m,n)
\]
rational. The prime tails $\{5,7,11\}$ and
$\{5,7,11,13,17,19\}$ satisfy the hypotheses of
Lemma~\ref{lem:shifts}.
The checks verify \eqref{eq:tail-norm}--\eqref{eq:tail-residual}
exactly. For the first tail, for example,
$s_E=167/385$ and
$q_K[h_E-f]=148819/27889$.
This residual is not small; the finite check tests an identity,
not proximity to the asymptotic regime.

For the invalid tail $\{2,3,5\}$, the norm and mixed-form identities
both fail. The support primes overlap the tail, and the shifted supports
can collide at index $6$. This control tests the actual missing
hypothesis. The verification receipt also records an earlier ineffective
negative-control choice; it is not reported as a successful control.

For $L=1$, write $H_j=\sum_{k\le j}1/k$, with $H_0=0$.
Formula~\eqref{eq:finite-gram} becomes
\[
 \ip{e_m}{A_Ne_n}
 =\frac{\gcd(m,n)}{\sqrt{mn}}\,
   \frac{H_{\lfloor N/\lcm(m,n)\rfloor}}{H_N}.
\]
Independent rowwise and pairwise evaluations agree exactly for
$N=12,31,100$. At $N=12$, the energy of the displayed finite vector
is $374167/172042$.
These checks concern finite identities only. They do not evaluate
resolvents, establish a uniform convergence theorem, or certify an
infinite prime tail. All limiting conclusions above follow from proofs.

\paragraph{Limits of the conclusion.}
Finite positivity alone neither proves nor disproves closability.
For example, when $\sum_p r_p^2<\infty$,
Proposition~\ref{prop:singularity} gives a positive closable form,
and its divergent-tail argument is unavailable.
In Theorem~\ref{thm:main}, the summability branch must likewise be
decided rather than inferred from the critical power $k^{-1/2}$ alone.

The divergent-case obstruction concerns nonnegative closed forms
agreeing with $K$ on $\cc$ inside the specified unweighted $\ell^2$ space.
It does not exclude changes of Hilbert space, other normalizations,
indefinite constructions, or unrelated spectral representations.
No target Euler factors, root numbers, zero correspondence, or
Hilbert--P\'olya realization follow from these arithmetic Gram limits.
